\documentclass[11pt, a4paper]{article}
\usepackage[utf8]{inputenc}
\usepackage[ngerman]{babel}
\usepackage[T1]{fontenc}
\usepackage{geometry}
\geometry{a4paper, left=20mm, right=20mm, top=20mm, bottom=20mm}
\usepackage{helvet}
\renewcommand{\familydefault}{\sfdefault}
\usepackage{enumitem}
\usepackage{titlesec}
\usepackage{hyperref}
\usepackage{xcolor}
\usepackage{tcolorbox}
\usepackage{amssymb}
\usepackage{longtable}
\usepackage{array}
\usepackage{booktabs}
\usepackage{fancyvrb}
\usepackage{listings}

\lstset{basicstyle=\ttfamily\small, breaklines=true, frame=single, backgroundcolor=\color{gray!10}}

\titleformat{\section}{\large\bfseries}{}{0em}{}
\titleformat{\subsection}{\normalsize\bfseries}{}{0em}{}
\titleformat{\subsubsection}{\small\bfseries}{}{0em}{}

\setlength{\parindent}{0pt}
\setlength{\parskip}{0.6em}

\definecolor{noteblue}{RGB}{0,90,160}
\definecolor{notebg}{RGB}{230,243,255}
\definecolor{linkcolor}{RGB}{0,70,140}

\hypersetup{colorlinks=true, linkcolor=linkcolor, urlcolor=linkcolor}

\begin{document}

\begin{center}
\textbf{\LARGE Arbeitsexpos\'e}\\[0.3em]
\textbf{\Large 3D Gaussian Splatting (3DGS) Scan-to-Building Information Modeling (BIM) Pipeline}\\[1em]
\textit{Arbeitsgrundlage f\"ur das formale Expos\'e an den Dozenten}\\[0.3em]
\small Stand: \today
\end{center}

\begin{tcolorbox}[colback=notebg, colframe=noteblue, title=\textbf{Zweck dieses Dokuments}]
Enth\"alt alle technischen Details, Repositories, Paper-Referenzen, Installations-TODOs und den Pipeline-Aufbau. Dient als Vorbereitung f\"ur das formale Kurzexpos\'e.
\end{tcolorbox}

\begin{tcolorbox}[colback=green!5, colframe=green!50!black, title=\textbf{Verf\"ugbare Hardware}]
\textbf{Graphics Processing Unit (GPU):} NVIDIA RTX 4000 Ada Generation -- \textbf{20\,GB Video Random Access Memory (VRAM)}, 6144 CUDA Cores, 160-bit GDDR6\\
\textbf{RAM:} 64\,GB\\
\textbf{Bewertung:} 20\,GB VRAM sind ausreichend f\"ur alle Pipeline-Komponenten (SuGaR: 12\,GB Min., Gaussian Grouping: 8--12\,GB, gsplat: 8\,GB+). Kein Cloud-Rechner notwendig.
\end{tcolorbox}

% ============================================================
\section{1. Titel und Forschungsfrage}

\textbf{Arbeitstitel:}\\
\textit{3D Gaussian Splatting vs. klassische Photogrammetrie: Geometrische Validierung, Segmentierung und Centerline-Extraktion f\"ur Leitungsinfrastruktur (Scan-to-BIM)}

\textbf{Zentrale Forschungsfrage (BA):}\\
\textit{L\"asst sich aus einem 3DGS-Modell \"uber SuGaR-Mesh-Extraktion und Centerline-Algorithmen eine geometrisch verwertbare 3D-Achslinie f\"ur r\"ohrenf\"ormige Infrastruktur ableiten, und wie genau ist diese im Vergleich zur Real-Time Kinematic Global Navigation Satellite System (RTK-GNSS)-Referenz?}

\begin{tcolorbox}[colback=blue!5, colframe=blue!50!black, title=\textbf{Kernbeitrag (Der Automatisierungs-Vorteil)}]
Eigentlich ist der gro\ss e Vorteil bei diesem Ansatz, dass sich lineare Objekte nicht nur in Vektorform erzeugen lassen, sondern auch als segmentierte Meshes direkt aus der realen Welt, die von vornherein georeferenziert sind. \\
Das ist der \textbf{Riesen-Vorteil} an dieser Arbeit: Die Pipeline wird so zusammengef\"uhrt, dass in der Praxis nur noch ein Drohnenflug und wenige Ground Control Points (GCPs) ben\"otigt werden. Am Ende liefert die Pipeline vollautomatisch ein segmentiertes 3D-Objekt. Da dieses \"uber 3D Gaussian Splatting und die SuGaR-Extraktion l\"auft, passt sich das Mesh physikalisch korrekt an die Oberfl\"ache an und sieht geometrisch deutlich besser (und lochfreier) aus, als wenn es rein \"uber klassische Photogrammetrie (Multi-View Stereo, MVS-Punktwolken-Triangulierung) erstellt w\"urde. Dadurch kann im letzten Schritt eine saubere Geoinformationssystem (GIS)-Centerline extrahiert werden.
\end{tcolorbox}

% ============================================================
\section{2. Pipeline-Architektur}

\begin{lstlisting}[frame=none, backgroundcolor=\color{white}]
Kamera-Video --> COLMAP (Structure-from-Motion, SfM) --> 3DGS (gsplat) --> Georef. (Helmert 7P)
                                                        |
               Gaussian Grouping (Segm.) --> SuGaR (Mesh) --> Centerline (DGtal)
                                                                    |
                                                              GIS-Import
\end{lstlisting}

\textbf{Umsetzung:} Die gesamte Pipeline wird durch ein zentrales Python-Orchestrierungs-Skript gesteuert. Dieses \"ubernimmt das automatische Wechseln der Conda-Umgebungen (via \texttt{subprocess}), das Fehler-Handling sowie das Time-Tracking f\"ur die sp\"atere Performance-Evaluation.

\textbf{Schritte im Detail:}

\begin{center}
\small
\begin{tabular}{|c|p{2.8cm}|p{3cm}|p{3.5cm}|p{3.5cm}|}
\hline
\textbf{\#} & \textbf{Schritt} & \textbf{Tool} & \textbf{Input} & \textbf{Output} \\
\hline
1 & Bilderfassung & Kamera (Video) & Campus-Testfeld & Bildsequenz (JPEG) \\
\hline
2 & SfM & COLMAP & Bilder & Sparse Point Cloud (PC) + Kameraposen \\
\hline
3 & 3DGS Training & gsplat & Sparse PC + Kameras & 3DGS-Modell (.ply) \\
\hline
4 & Georeferenzierung & Eigenes Skript & GCP-Koord. + Modell & Georef. 3DGS \\
\hline
5 & Segmentierung & \textit{Vergleichsstudie} & 3DGS + Annotationen & Segm. Gaussians \\
\hline
6 & Mesh-Extraktion & SuGaR & Segm. Gaussians & 3D-Mesh (.obj/.ply) \\
\hline
7 & Centerline & DGtal / Python & Kabel-Mesh & 3D-Linie (Polyline) \\
\hline
8 & Geometrische H\"ohenkorrektur & Python & Centerline + Radius & Korrigierte Centerline (Z-Shift) \\
\hline
9 & GIS-Import & ArcGIS / QGIS & Georef. Linie & Shapefile / GeoJSON \\
\hline
10 & Evaluation & Python (numpy) & Centerline vs. RTK & Root Mean Square Error (RMSE) [cm] \\
\hline
\end{tabular}
\end{center}

% ============================================================
\section{3. Segmentierungsvergleich (neu)}

F\"ur die Isolation der Kabel-Gaussians werden \textbf{drei Methoden} systematisch verglichen. Dies ergibt eine zus\"atzliche, methodisch wertvolle Forschungsfrage.

\begin{center}
\small
\begin{tabular}{|p{2.5cm}|p{2.5cm}|p{3cm}|p{2cm}|p{2.5cm}|}
\hline
\textbf{Methode} & \textbf{Typ} & \textbf{Prinzip} & \textbf{Training?} & \textbf{Paper} \\
\hline
(A) Farb-Filterung & Baseline & Schwellwert auf Red-Green-Blue (RGB)/Hue-Saturation-Value (HSV)-Attribute der Splats & Nein & -- (eigene Impl.) \\
\hline
(B) FlashSplat & Optimierung & 2--3 manuelle 2D-Annotationen $\rightarrow$ optimale 3D-Maske & Nein & Shen et al. (ECCV 2024) \\
\hline
(C) Gaussian Grouping & Training & Identity Encoding via DEVA/SAM-Masken & Ja & Ye et al. (ECCV 2024) \\
\hline
\end{tabular}
\end{center}

\textbf{Bewertungskriterien:}
\begin{itemize}[noitemsep]
    \item \textbf{Segmentierungsqualit\"at:} Intersection over Union (IoU) / Precision / Recall der isolierten Kabel-Gaussians
    \item \textbf{Aufwand:} Setup-Komplexit\"at, Annotationszeit, Dependencies
    \item \textbf{Auswirkung auf Centerline:} RMSE-Differenz der resultierenden Achslinie
\end{itemize}

\begin{tcolorbox}[colback=notebg, colframe=noteblue]
\textbf{Wissenschaftlicher Mehrwert:} Der dreistufige Vergleich (manuell $\rightarrow$ trainingsfrei $\rightarrow$ trainingsbasiert) zeigt systematisch, ab welcher Komplexit\"atsstufe sich der Aufwand f\"ur die Kabelisolation lohnt. Gleichzeitig dient die Farbfilterung als Baseline, gegen die die ML-Methoden quantitativ bewertet werden.
\end{tcolorbox}

\textbf{Referenz-Survey:} He et al. (2025): \textit{``A Survey on 3D Gaussian Splatting Applications: Segmentation, Editing, and Generation.''} ArXiv: 2508.09977. Umfassendster Survey zu 3DGS-Segmentierung mit 50+ Methoden -- dient als Grundlage f\"ur die Methodenauswahl und das State-of-the-Art-Kapitel.

% ============================================================
\section{4. Scope-Trennung PA / BA}

\subsection{Projektarbeit (PA) -- 7 European Credit Transfer System (ECTS), 4 Wochen netto}
\textbf{Scope:} Testfeldaufbau, Datenerfassung, Proof of Concept (Schritte 1--3).

\begin{center}
\small
\begin{tabular}{|p{3.5cm}|p{6.5cm}|p{2.5cm}|}
\hline
\textbf{TODO} & \textbf{Details} & \textbf{Zeit} \\
\hline
Testfeld aufbauen & Gartenschlauch + Schn\"ure auf Campus & 1 Tag \\
\hline
GCPs einmessen & 8--12 GCPs mit RTK-GNSS (HS-Ger\"ate) & 1--2 Tage \\
\hline
TLS-Referenzscan & Laserscanner-Aufnahme des Testfelds & 0.5--1 Tag \\
\hline
Videoaufnahme & Kamerasequenz (Nadir + Oblique) & 0.5 Tage \\
\hline
COLMAP installieren & Pre-built Windows Binary & 0.5 Tage \\
\hline
COLMAP Pipeline & Feature Extr. $\rightarrow$ Matching $\rightarrow$ Sparse & 1 Tag \\
\hline
MVS-Baseline & Dense Reconstruction (Metashape/COLMAP) & 1--2 Tage \\
\hline
gsplat installieren & Conda + pip install gsplat & 1 Tag \\
\hline
3DGS trainieren & gsplat auf COLMAP-Output & 1--2 Tage \\
\hline
Visueller Vergleich & Screenshots MVS vs. 3DGS & 1 Tag \\
\hline
Rechenzeit-Benchmark & Tabelle: MVS vs. 3DGS & 0.5 Tage \\
\hline
PA schreiben & 15--25 Seiten & 5--7 Tage \\
\hline
\end{tabular}
\end{center}

\textbf{Forschungsfragen PA:}
\begin{enumerate}[noitemsep]
    \item Lassen sich texturlose und feine Strukturen durch 3DGS visuell vollst\"andiger rekonstruieren als durch MVS?
    \item Wie verhalten sich die Rechenzeiten der dichten Rekonstruktion?
\end{enumerate}

\textbf{Ergebnis PA:} Validierter Datensatz, lauff\"ahige Pipelines, erster Vergleich.

\subsection{Bachelorarbeit (BA) -- 15 ECTS, 10 Wochen netto}
\textbf{Scope:} Georeferenzierung, \textbf{Segmentierungsvergleich (3 Methoden)}, Mesh, Centerline, Evaluation (Schritte 4--9).

\begin{center}
\small
\begin{tabular}{|p{3.5cm}|p{6.5cm}|p{2.5cm}|}
\hline
\textbf{TODO} & \textbf{Details} & \textbf{Zeit} \\
\hline
Georef-Skript & Helmert 7P (numpy/scipy) & 2--3 Tage \\
\hline
(A) Farb-Filterung & HSV-Schwellwert auf Splat-Attribute & 1--2 Tage \\
\hline
(B) FlashSplat & 2--3 Annotationen $\rightarrow$ 3D-Maske & 2--3 Tage \\
\hline
(C) Gaussian Grouping & DEVA-Masken + Identity Encoding & 3--5 Tage \\
\hline
SuGaR installieren & Conda (Python 3.9, CUDA 11.8) & 1--2 Tage \\
\hline
SuGaR Windows-Fix & Pfadtrenner mit pathlib fixen & 0.5 Tage \\
\hline
SuGaR Mesh & Poisson-Rekonstruktion Kabel-Gaussians & 1--2 Tage \\
\hline
Centerline-Algorithmus & DGtal oder Python-Reimplementierung & 3--5 Tage \\
\hline
GIS-Import & Export Shapefile $\rightarrow$ ArcGIS/QGIS & 1--2 Tage \\
\hline
RMSE-Evaluation & Centerline vs. RTK-Referenzlinie & 2--3 Tage \\
\hline
Wirtschaftlichkeit & Kosten/Zeit vs. klass. Vermessung & 2--3 Tage \\
\hline
BA schreiben & 40--60 Seiten & 15--20 Tage \\
\hline
\end{tabular}
\end{center}

\textbf{Forschungsfragen BA:}
\begin{enumerate}[noitemsep]
    \item Wie genau ist die extrahierte Centerline im Vergleich zur RTK-Referenz (RMSE)?
    \item \textbf{Welche der drei Segmentierungsmethoden (Farb-Filter, FlashSplat, Gaussian Grouping) liefert die robusteste Isolation bei geringstem Aufwand?}
    \item \textbf{Praxis-Problem:} Wie l\"asst sich der systematische H\"ohenfehler der Centerline algorithmisch korrigieren (Z-Shift), der durch die unvollst\"andige optische Erfassung des Kabels im Graben (Halbzylinder-Geometrie) entsteht?
    \item Ab welchem Objektdurchmesser scheitert die Pipeline?
    \item Ist 3DGS eine wirtschaftlichere Alternative f\"ur die Bau\"uberwachung?
\end{enumerate}

% ============================================================
\section{5. GitHub-Repositories}

\begin{center}
\small
\begin{tabular}{|p{3cm}|p{8cm}|p{2cm}|}
\hline
\textbf{Tool} & \textbf{URL} & \textbf{Lizenz} \\
\hline
COLMAP & \url{https://github.com/colmap/colmap} & BSD-3 \\
\hline
gsplat & \url{https://github.com/nerfstudio-project/gsplat} & Apache-2.0 \\
\hline
SuGaR & \url{https://github.com/Anttwo/SuGaR} & Research \\
\hline
Gaussian Grouping & \url{https://github.com/lkeab/gaussian-grouping} & Apache-2.0 \\
\hline
\textbf{FlashSplat} & \url{https://github.com/florinshen/FlashSplat} & Research \\
\hline
DEVA & \url{https://github.com/hkchengrex/Tracking-Anything-with-DEVA} & -- \\
\hline
DGtal & \url{https://github.com/DGtal-team/DGtal} & LGPL \\
\hline
Original 3DGS & \url{https://github.com/graphdeco-inria/gaussian-splatting} & Research \\
\hline
\end{tabular}
\end{center}

% ============================================================
\section{6. Paper-Referenzen}

\subsection{Kern-Paper (direkt in Pipeline genutzt)}
\begin{enumerate}[noitemsep]
    \item \textbf{3D Gaussian Splatting} -- Kerbl et al. (SIGGRAPH 2023)\\
    \textit{``3D Gaussian Splatting for Real-Time Radiance Field Rendering.''}\\
    ACM Transactions on Graphics 42(4). Best Paper SIGGRAPH 2023.

    \item \textbf{SuGaR} -- Gu\'edon, Lepetit (CVPR 2024)\\
    \textit{``SuGaR: Surface-Aligned Gaussian Splatting for Efficient 3D Mesh Reconstruction.''}\\
    ArXiv: 2311.12775

    \item \textbf{Gaussian Grouping} -- Ye et al. (ECCV 2024)\\
    \textit{``Gaussian Grouping: Segment and Edit Anything in 3D Scenes.''}

    \item \textbf{FlashSplat} -- Shen et al. (ECCV 2024)\\
    \textit{``FlashSplat: 2D to 3D Gaussian Splatting Segmentation Solved Optimally.''}\\
    Trainingsfreie Segmentierung durch Optimierung auf 2D-Annotationen.

    \item \textbf{Centerline-Extraktion} -- Kerautret et al. (DGCI 2016)\\
    \textit{``3D Geometric Analysis of Tubular Objects based on Surface Normal Accumulation.''}

    \item \textbf{COLMAP} -- Sch\"onberger, Frahm (CVPR 2016)\\
    \textit{``Structure-from-Motion Revisited.''}
\end{enumerate}

\subsection{Vergleichsstudien und Benchmarks}
\begin{enumerate}[noitemsep, start=6]
    \item \textbf{GauU-Scene V2 Benchmark} -- Inverse Korrelation PSNR vs. geometrische Genauigkeit
    \item \textbf{ISPRS Annals 2025} -- Geod\"atisch relevanteste Vergleichsstudie 3DGS vs. Photogrammetrie
    \item \textbf{GeoRefGS} -- \textit{``Towards Georeferenced 3DGS from UAV Platforms''} (MDPI 2025)\\
    Erw\"ahnt als State-of-the-Art, nicht genutzt (kein \"offentlicher Code)
    \item \textbf{3DGS Applications Survey} -- He et al. (2025)\\
    \textit{``A Survey on 3D Gaussian Splatting Applications: Segmentation, Editing, and Generation.''}\\
    ArXiv: 2508.09977. Umfassendster Survey mit 50+ Segmentierungsmethoden.
\end{enumerate}

\subsection{Varianten (State-of-the-Art)}
\begin{enumerate}[noitemsep, start=9]
    \item \textbf{Mip-Splatting} (CVPR 2024) -- Anti-Aliasing f\"ur Multi-Scale
    \item \textbf{2D Gaussian Splatting} (SIGGRAPH 2024) -- Planare Gaussians
    \item \textbf{GaussianPro} (ICML 2024) -- Texturlose Regionen
    \item \textbf{Scaffold-GS} (CVPR 2024) -- Speicher-Reduktion
    \item \textbf{DroneSplat} -- Drohnenspezifische 3DGS-Pipeline
\end{enumerate}

\subsection{Large-Scale 3DGS (Skalierbarkeit)}
\begin{enumerate}[noitemsep, start=14]
    \item \textbf{VastGaussian} -- Lin et al. (CVPR 2024)\\
    \textit{``VastGaussian: Vast 3D Gaussians for Large Scene Reconstruction.''}\\
    Automatische Partitionierung gro\ss er Szenen in Tiles mit seamless Merging.
    \item \textbf{CityGaussian} -- Liu et al. (2024)\\
    \textit{``CityGaussian: Real-time High-quality Large-Scale Scene Rendering.''}\\
    Level-of-Detail f\"ur stadtweite 3DGS-Rekonstruktionen.
\end{enumerate}

% ============================================================
\section{7. Installations-Checkliste}

\subsection{Voraussetzung: Hardware}
\begin{itemize}[noitemsep]
    \item[$\checkmark$] \textbf{NVIDIA RTX 4000 Ada -- 20\,GB VRAM} (vorhanden)
    \item[$\checkmark$] \textbf{64\,GB RAM} (vorhanden)
    \item[$\square$] NVIDIA Treiber $\geq$ 525.x (CUDA 11.3--11.8)
    \item[$\square$] Conda/Miniconda installiert
    \item[$\square$] Git installiert
    \item[$\square$] Visual Studio Build Tools (C++ Compiler)
\end{itemize}

\subsection{COLMAP}
\begin{itemize}[noitemsep]
    \item[$\square$] Pre-built Binary von \url{https://demuc.de/colmap/} (CUDA-Version)
    \item[$\square$] Entpacken und PATH setzen
    \item[$\square$] Test: \texttt{colmap -h}
\end{itemize}

\subsection{gsplat (3DGS Training)}
\begin{lstlisting}
conda create -n gsplat python=3.10 -y
conda activate gsplat
pip install torch torchvision --index-url https://download.pytorch.org/whl/cu118
pip install gsplat
\end{lstlisting}

\subsection{SuGaR (Mesh-Extraktion)}
\begin{lstlisting}
git clone https://github.com/Anttwo/SuGaR.git --recursive
cd SuGaR
python install.py   # Erstellt conda env "sugar"
conda activate sugar
\end{lstlisting}
\begin{itemize}[noitemsep]
    \item[$\square$] Windows-Pfade fixen: \texttt{/} $\rightarrow$ \texttt{os.path.join()} in relevanten Skripten
\end{itemize}

\subsection{FlashSplat (Segmentierung -- trainingsfrei)}
\begin{lstlisting}
git clone https://github.com/florinshen/FlashSplat.git
cd FlashSplat
pip install -r requirements.txt
\end{lstlisting}

\subsection{Gaussian Grouping (Segmentierung -- trainingsbasiert)}
\begin{lstlisting}
git clone https://github.com/lkeab/gaussian-grouping.git
cd gaussian-grouping
conda create -n gaussian_grouping python=3.8 -y
conda activate gaussian_grouping
conda install pytorch==1.12.1 torchvision==0.13.1 cudatoolkit=11.3 -c pytorch
pip install plyfile==0.8.1 tqdm scipy wandb opencv-python scikit-learn lpips
pip install submodules/diff-gaussian-rasterization
pip install submodules/simple-knn
\end{lstlisting}

\subsection{DGtal (Centerline)}
\begin{lstlisting}
git clone https://github.com/DGtal-team/DGtal.git
cd DGtal && mkdir build && cd build
cmake .. && make
\end{lstlisting}
Alternativ: Centerline-Algorithmus (Normalenakkumulation) in Python mit Open3D reimplementieren.

\subsection{Georeferenzierung (eigenes Skript)}
\begin{lstlisting}
pip install numpy scipy   # Kein extra Environment noetig
\end{lstlisting}
\begin{itemize}[noitemsep]
    \item[$\square$] CSV: \texttt{ID, X\_lokal, Y\_lokal, Z\_lokal, X\_welt, Y\_welt, Z\_welt}
    \item[$\square$] Helmert 7P mit \texttt{scipy.spatial.transform}
    \item[$\square$] Transformation auf alle Gaussian-Positionen anwenden
\end{itemize}

% ============================================================
\section{8. Evaluation-Framework (D1--D6)}

\begin{center}
\small
\begin{tabular}{|p{2cm}|p{4cm}|p{3cm}|c|c|}
\hline
\textbf{Dim.} & \textbf{Metrik} & \textbf{Tool} & \textbf{PA} & \textbf{BA} \\
\hline
D1: Geometrie & RMSE horiz./vert. [cm] vs. TLS & CloudCompare / numpy & -- & $\checkmark$ \\
\hline
D2: Visuell & PSNR, SSIM, LPIPS & gsplat metrics & $\checkmark$ & $\checkmark$ \\
\hline
D3: Manuell & Visueller Vergleich & Screenshots & $\checkmark$ & $\checkmark$ \\
\hline
D4: ML-Erkennung & Precision, Recall, F1 & sklearn & -- & $\checkmark$ \\
\hline
D5: Kosten & Zeit, Kosten, Hardware & Tabelle & $\checkmark$ & $\checkmark$ \\
\hline
D6: Tile-Konsistenz & Overlap-RMSE [cm] & numpy & -- & $\checkmark$ \\
\hline
\end{tabular}
\end{center}

\begin{tcolorbox}[colback=green!5, colframe=green!50!black, title=\textbf{Referenzdaten (verf\"ugbar an der Hochschule)}]
\textbf{RTK-GNSS:} F\"ur GCP-Einmessung und Georeferenzierung (Genauigkeit: $<$\,1\,cm).\\
\textbf{Terrestrischer Laserscanner (TLS):} F\"ur dichte Referenz-Punktwolke des gesamten Testfelds. Daraus wird eine \textbf{TLS-Referenz-Centerline} extrahiert, gegen die die 3DGS-Centerline evaluiert wird (Linie-zu-Linie-Vergleich, nicht nur Punkt-zu-Punkt).\\
$\rightarrow$ \textbf{Doppelte Referenz:} GCPs f\"ur Georeferenzierung, TLS f\"ur geometrische Evaluation.
\end{tcolorbox}

% ============================================================
\section{9. Testfeld-Design (Campus, ca. 20--30\,m)}

Die Kamera wird auf ca. \textbf{2\,m H\"ohe} gef\"uhrt; bei einer realen Drohnenh\"ohe von 10\,m ergibt sich ein \textbf{Skalierungsfaktor von 5$\times$}. Die Testobjekt-Durchmesser erzeugen somit bei 2\,m die gleiche Pixelabdeckung (GSD) wie reale Kabel bei 10\,m Flugh\"ohe:

\begin{center}
\small
\begin{tabular}{|p{3cm}|p{2cm}|p{3.5cm}|p{4.5cm}|}
\hline
\textbf{Testobjekt} & \textbf{$\varnothing$ Campus} & \textbf{Simuliert} & \textbf{$\varnothing$ real (10\,m H\"ohe)} \\
\hline
Gartenschlauch (dick) & $\sim$40\,mm & DC-Erdkabel & 20\,cm (200\,mm) \\
\hline
Gartenschlauch (d\"unn) & $\sim$16\,mm & Steuerkabel & 8\,cm (80\,mm) \\
\hline
Paketschnur / Seil & $\sim$10\,mm & Lichtwellenleiter & 5\,cm (50\,mm) \\
\hline
Nylonschnur & $\sim$3\,mm & Grenztest & $<$\,2\,cm \\
\hline
GCP-Targets & -- & Passpunkte & 8--12 St\"uck, RTK-GNSS ($<$\,1\,cm) \\
\hline
TLS-Scan & -- & Referenz & Terrestr. Laserscanner (HS-Ger\"at) \\
\hline
\end{tabular}
\end{center}

% ============================================================
\section{10. Skalierbarkeit: Tile-basiertes Processing}

Standard-3DGS ist f\"ur kleine bis mittlere Szenen optimiert ($\sim$50\,m$^2$). Reale Kabeltrassen erstrecken sich \"uber \textbf{100\,m bis 1\,km}. Die L\"osung ist ein \textbf{georeferenziertes Tiling}, das in der Photogrammetrie Standard ist.

\subsection{Tile-Strategie}

\begin{lstlisting}[frame=none, backgroundcolor=\color{white}]
1km Trasse --> Aufteilen in N Tiles a 50m (mit 10m Ueberlappung)
  |
  Pro Tile: COLMAP --> 3DGS --> Georef --> Segm. --> SuGaR --> Centerline
  |
  Alle Centerlines in georef. Koordinaten (UTM)
  |
  GIS-Merge: Spatial Join auf Ueberlappungszonen --> durchgehende Linie
\end{lstlisting}

\begin{center}
\small
\begin{tabular}{|p{3cm}|p{2.5cm}|p{3cm}|p{4cm}|}
\hline
\textbf{Parameter} & \textbf{Empfehlung} & \textbf{Begr\"undung} & \textbf{Auswirkung} \\
\hline
Tile-L\"ange & 50\,m & Optimale 3DGS-Szenengr\"o\ss e & $\sim$500k Gaussians/Tile \\
\hline
\"Uberlappung & 10\,m (20\%) & Genug Kontext f\"ur COLMAP & Seamless Centerline-Join \\
\hline
GCPs pro Tile & 3--4 Minimum & F\"ur Helmert-Transformation & $<$\,2\,cm Georef.-Fehler \\
\hline
\end{tabular}
\end{center}

\subsection{Merge-Strategie im GIS}
\begin{enumerate}[noitemsep]
    \item Alle Tile-Centerlines als Polylines in ein GIS laden (gleiche Projektion)
    \item In der \"Uberlappungszone: \textbf{gewichtete Mittelung} der Positionen
    \item \textbf{Spatial Join} auf $<$\,0{,}5\,m Abstand $\rightarrow$ durchgehende Linie
    \item \textbf{Overlap-RMSE} berechnen: Abweichung der Centerlines beider Tiles in der \"Uberlappung
\end{enumerate}

\begin{tcolorbox}[colback=notebg, colframe=noteblue]
\textbf{Wissenschaftlicher Mehrwert:} Der Overlap-RMSE ist eine \textbf{neue Evaluationsdimension} (D6), die die Konsistenz der Pipeline \"uber Tile-Grenzen hinweg misst. F\"ur das Campus-Testfeld (5--10\,m) ist Tiling nicht notwendig -- es kann aber als Proof-of-Concept getestet werden, indem man die Szene k\"unstlich in 2--3 Tiles aufteilt.
\end{tcolorbox}

\textbf{Referenzen:} VastGaussian (CVPR 2024) und CityGaussian (2024) zeigen, dass 3DGS auf Stadtma\ss stab skaliert werden kann. Der hier gew\"ahlte manuelle Tiling-Ansatz ist f\"ur linienf\"ormige Infrastruktur \textbf{einfacher und robuster}, da die Georeferenzierung die r\"aumliche Konsistenz sicherstellt.

% ============================================================
\section{11. Risikomatrix}

\begin{center}
\small
\begin{tabular}{|p{3.5cm}|p{2cm}|p{1.5cm}|p{5.5cm}|}
\hline
\textbf{Risiko} & \textbf{Wahrsch.} & \textbf{Impact} & \textbf{Mitigation} \\
\hline
SuGaR Windows-Pfade & Hoch & Niedrig & Trivial fixbar mit pathlib \\
\hline
Kabel nicht erkannt & Mittel & Hoch & Fallback: manuelle Farb-Filterung \\
\hline
Centerline schlecht & Mittel & Hoch & Alternative: Open3D Skeleton \\
\hline
VRAM $<$ 12\,GB & \textit{Entf\"allt} & -- & RTX 4000 Ada mit 20\,GB vorhanden \\
\hline
RTK-GNSS / TLS nicht verf. & \textit{Gering} & Mittel & Ger\"ate an HS vorhanden, fr\"uhzeitig reservieren \\
\hline
Mesh verrauscht & Mittel & Mittel & SuGaR \texttt{dn\_consistency} nutzen \\
\hline
Szene zu gro\ss{} & Mittel & Mittel & Tile-basiertes Processing (50\,m Tiles) \\
\hline
\end{tabular}
\end{center}

\subsection{Externe Machbarkeitsanalyse (ValyuAI)}
Eine unabh\"angige Machbarkeitspr\"ufung best\"atigt f\"ur dieses Forschungsdesign die Statusfreigabe \textbf{START GR\"UN}. 
\begin{itemize}[noitemsep]
    \item \textbf{Machbarkeit:} Sehr hoch (alle Komponenten verifiziert, Hardware mit RTX 4000 20\,GB ausreichend).
    \item \textbf{Wissenschaftlicher Mehrwert:} Sehr hoch. Es schlie\ss t eine echte Forschungsl\"ucke (3DGS-Genauigkeit f\"ur Centerline-Extraktion).
    \item \textbf{Fallbacks (Risiko-Minimierung):} Falls \textit{Gaussian Grouping} zu rechenintensiv wird, kann auf \textit{FlashSplat} reduziert werden. Falls \textit{DGtal} h\"angt, kann das Mesh exportiert und extern (z.B. Open3D) verarbeitet werden.
    \item \textbf{Publikationspotenzial:} Es sind bis zu 3 Papers ableitbar (Segmentierungs-Vergleich, Scan-to-BIM Genauigkeitsanalyse gegen TLS, Open-Source Toolkit Report).
\end{itemize}

% ============================================================
\section{12. Diskussionspunkte f\"ur den Dozenten}

\begin{enumerate}[noitemsep]
    \item \textbf{Scope PA vs. BA:} PA = Datenerhebung + Vergleich (MVS vs. 3DGS). BA = Scan-to-BIM Pipeline.
    \item \textbf{Segmentierung:} Dreistufiger Vergleich: Farb-Filterung vs. FlashSplat vs. Gaussian Grouping.
    \item \textbf{Georeferenzierung:} Klassische Helmert-Transformation \"uber COLMAP GCPs.
    \item \textbf{Skalierbarkeit:} Tile-basiertes Processing (50\,m Tiles, 10\,m Overlap) f\"ur Trassen bis 1\,km. Machbarkeitstest auf Campus durch k\"unstliches Aufteilen in 2--3 Tiles.
    \item \textbf{Referenzdaten:} RTK-GNSS (GCPs) + TLS-Scan (Referenz-Centerline) -- beides an der HS verf\"ugbar.
    \item \textbf{Testfeld:} Campus ausreichend, oder Drohnenflug?
    \item \textbf{KI-Nutzung:} Dokumentation gem\"a\ss{} PABA-Hinweise Kap.~3.5.
    \item \textbf{Hardware:} RTX 4000 Ada (20\,GB) + 64\,GB RAM vorhanden.
    \item \textbf{Centerline:} DGtal (C++) oder eigene Python-Reimplementierung?
    \item \textbf{Survey:} He et al. (2025) als Grundlage f\"ur SotA-Kapitel (50+ Methoden).
\end{enumerate}

\end{document}
